\documentclass[10pt,twocolumn]{article}
\usepackage[margin=0.7in]{geometry}
\usepackage{amsmath,booktabs,graphicx}
\title{QuoteGate: Information-Driven Trading Under Latency and Execution Constraints}
\author{QuoteGate team (HackRU 2026)}
\date{}
\begin{document}
\maketitle

\section{Problem and contribution}
% TODO(results): one paragraph. Systems + quant contributions; synthetic setting stated up front.
We ask, under a declared synthetic market model, whether a displayed-quantity imbalance signal from
one reference venue carries predictive information, and whether any advantage survives spread,
fees, execution delay, and venue failures. All data are synthetic; nothing here is evidence about
real markets.

\section{System}
% TODO(figure): small architecture diagram (market -> venues -> aggregator -> strategy -> router -> venue execution -> ledger).
Venue simulators stream quotes over persistent TCP connections to an aggregator that tracks
per-venue sessions, sequence order, and local-receipt freshness. The strategy observes only the
reference venue's accepted quotes and freshness, plus its own confirmed position and order state.

\section{Signal and policy}
For bid and ask displayed quantities $b_t, a_t \ge 1$ at the reference venue,
\begin{equation}
  I_t = \frac{b_t - a_t}{b_t + a_t} \in (-1, 1).
\end{equation}
With threshold $\theta$ and limit $q_{\max}$, the target position is
\begin{equation}
  q^{*}_t = \begin{cases} +q_{\max} & I_t > \theta \\ -q_{\max} & I_t < -\theta \\ 0 & \text{otherwise,} \end{cases}
\end{equation}
and the order is $\Delta q_t = q^{*}_t - q_t$ (confirmed position), clipped to the order limit.
No order is sent while the signal is stale, an order is outstanding, or an outcome is unknown.

\section{Synthetic market and execution}
Latent price $p_t$ in integer ticks; the reference venue quotes $[p_t, p_t + s]$, a second venue
quotes no better. With probability $m$ the price moves one tick, up with probability
$(1 + \beta I_t)/2$; $\beta = 0$ is the null regime. Orders are immediate-or-cancel limits at the
observed price, filled in full at the venue's then-current ask (bid) or rejected; acknowledgement
loss yields an \emph{unknown} outcome resolved only by reconciliation. Accounting:
$c \leftarrow c - p\,\Delta q - f$, $q \leftarrow q + \Delta q$.

\section{Experiments}
% TODO(results): fixed matrix, seeds, frozen parameters; results table; one plot.
\begin{table}[h]\centering\small
\begin{tabular}{lrrrr}
\toprule
Cell & Fills & Gross & Net & Hit rate \\
\midrule
% TODO(results)
\bottomrule
\end{tabular}
\caption{TODO(results).}
\end{table}

\section{Limitations}
% TODO(results): synthetic one-step signal, no queue/impact, IOC fill model, single instrument.

\begin{thebibliography}{9}\small
% NOTE: verify each entry against the primary source before submission.
\bibitem{cjp} \'A.~Cartea, S.~Jaimungal, J.~Penalva. \emph{Algorithmic and High-Frequency Trading}. Cambridge University Press, 2015.
\bibitem{gb} M.~D.~Gould, J.~Bonart. Queue imbalance as a one-tick-ahead price predictor in a limit order book. \emph{Market Microstructure and Liquidity}, 2016.
\bibitem{cks} R.~Cont, A.~Kukanov, S.~Stoikov. The price impact of order book events. \emph{Journal of Financial Econometrics}, 2014.
\end{thebibliography}
\end{document}
